% docs/manual/user/chapters/04-variables-references.tex
% 第 4 章：选择变量与参考数据集

\chapter{选择变量与参考数据集}

\openbench{} 的 reference 注册表内置近百个数据集，覆盖陆面模型评估的 9 大类核心变量。本章帮你回答两个问题：

\begin{itemize}
  \item \textbf{我能评估什么变量}（看变量分类）
  \item \textbf{这个变量用哪个参考数据集}（看变量到数据集映射）
\end{itemize}

完整数据集清单见附录 C，本章重点解释结构、用法、命令行工具。

\section{变量分类：9 大类}

\openbench{} 的变量按物理过程分为 9 大类：

\begin{itemize}
  \item \textbf{Bio}：生态系统结构（LAI、Biomass、Burned\_Area）
  \item \textbf{Carbon}：碳通量（GPP、NEE、NPP、Reco、Methane、Wetland\_Methane）
  \item \textbf{Water}：水通量与水储量（ET 系列、Runoff、Streamflow、Snow、TWS、Inundation、Precipitation）
  \item \textbf{Energy}：能量通量（Latent\_Heat、Sensible\_Heat、Net\_Radiation、Ground\_Heat、辐射四分量、Albedo）
  \item \textbf{Heat}：地表/土壤温度
  \item \textbf{Meteorology}：气象驱动量（Air\_Temperature、Tmax/Tmin、Humidity、Pressure、Wind、Cloud\_Cover）
  \item \textbf{Urban}：城市与农业（Crop、Urban、Water\_Use、Dam）
  \item \textbf{Composite}：合成或衍生指标
  \item \textbf{Other}：暂无固定分类
\end{itemize}

变量名严格遵守驼峰式 + 下划线（\var{Latent\_Heat} 而非 \texttt{LatentHeat} 或 \texttt{latent\_heat}）。常见标准变量名见附录 C 中数据集表的"变量"列。

\section{Grid 与 Station：两种数据形态}

\openbench{} 区分两种参考数据形态：

\begin{itemize}
  \item \textbf{grid}：规则经纬度网格。NetCDF 文件含 \texttt{lat}、\texttt{lon}、\texttt{time} 三个坐标。代表数据：\dataset{GLEAM\_v4.2a}（ET）、\dataset{ERA5LAND}（多变量）、\dataset{IMERG}（降水）
  \item \textbf{stn}（station）：离散站点观测。NetCDF 文件含 \texttt{site}（或类似的站点 id）+ \texttt{time}。代表数据：\dataset{FLUXNET\_PLUMBER2}（FLUXNET）、\dataset{GRDC\_Daily}（GRDC 流量站）、\dataset{Caravan\_Daily}
\end{itemize}

\subsection{4 种数据形态组合}

sim 与 ref 各 grid/stn → 4 种组合，\openbench{} 都支持：

\begin{longtable}{p{2cm} p{2cm} p{8.5cm}}
  \toprule
  sim & ref & 处理方式 \\
  \midrule
  \endhead
  grid & grid & 默认；两边重网格化到 \yamlkey{project.grid_res} 后逐像素对比 \\
  grid & stn & sim 在站点位置采样（最近邻 / 双线性），与 stn ref 对比；要 \yamlkey{fulllist} \\
  stn & grid & ref 在站点位置采样，与 stn sim 对比（罕见，但支持） \\
  stn & stn & 两边都按站点列表逐点对比；station\_list 优先取 ref/sim 的 \yamlkey{fulllist}，否则自动从 sim 目录扫描 \\
  \bottomrule
\end{longtable}

任一边是 stn 都触发"站点驱动模式"：评估按 station\_list 中每个站点独立做 metric。同站点 ID 在 sim 与 ref 中要能对上（缺失的站点会被跳过并 WARN）。

\subsection{站点数据的特殊配置}

如果你的 simulation 是 grid 模型而 reference 是 station，\openbench{} 会做"以站点位置抽样 grid"的处理。你需要在 \yamlkey{simulation} entry 里指定 \yamlkey{fulllist}（站点列表 CSV 路径）：

\begin{minted}{yaml}
simulation:
  CoLM2024_PLUMBER2:
    model: CoLM2024
    root_dir: /data/colm2024
    fulllist: /data/PLUMBER2/list/PLUMBER2.csv
\end{minted}

如果两边都是 stn（同站点列表的两个 dataset 直接对比），\yamlkey{fulllist} 仍可指定，但通常省略 —— \openbench{} 会自动扫描 sim 目录提取站点列表。详见第 5 章。

\subsection{min\_year\_threshold：站点最小数据长度过滤}

\yamlkey{project.min_year_threshold}（默认 \texttt{3}，单位：年）只对站点参考数据生效，是\emph{站点级}过滤器：

\begin{enumerate}
  \item \openbench{} 读站点 NC 后，对每个站逐一计算其有效数据时间跨度 = $\min(\text{站结束年}, \text{sim结束年}, \text{eyear}) - \max(\text{站开始年}, \text{sim开始年}, \text{syear})$
  \item 若该跨度 \texttt{<} \yamlkey{min_year_threshold}，整个站被剔除（不进入 station list）
  \item 通过的站才进入 evaluation 阶段
\end{enumerate}

实现位置：\modname{data.station_matcher} 与各 \modname{data.custom.*\_filter} 模块（如 GRDC、HydroWeb、ResOpsUS）。

\subsection*{挑值的指引}

\begin{itemize}
  \item \textbf{默认 3}：跨年间变率 metrics（KGE、相关性）所需的最低样本量经验值
  \item \textbf{设 1}：你正在做单年度评估或纯空间分析，不在意跨年统计稳健性
  \item \textbf{设 5--10}：发表论文中的 trend / 趋势 metric，希望排除短记录站点
  \item 注意 \emph{这个阈值与 \yamlkey{years} 一起作用}：缩小 \yamlkey{years} 区间会减少每个站的可用跨度，可能让原本通过的站被剔除
\end{itemize}

\textbf{陷阱}：\cli{check} 不会提前告诉你这个过滤器会剔除多少站；只在 run 阶段日志看到 \texttt{[INFO] Selected N stations} 与原始站数对比。如果选出的站数远低于预期，先把这个阈值调小再排查。

\section{分辨率后缀自动解析}

许多 reference 数据集提供多个分辨率版本，registry 用后缀区分：

\begin{itemize}
  \item \texttt{\_LowRes}：低分辨率（典型 0.5°）
  \item \texttt{\_MidRes}：中分辨率（典型 0.25°）
  \item 无后缀：原生分辨率，或单一可用版本
\end{itemize}

例如 \dataset{GLEAM\_v4.2a} 在 catalog 中实际是两个 entry：

\begin{itemize}
  \item \dataset{GLEAM\_v4.2a\_LowRes}（0.5°）
  \item \dataset{GLEAM\_v4.2a\_MidRes}（0.25°）
\end{itemize}

你在 \yamlkey{reference} 写时，可以：

\subsection{方式 1：写 base 名让 \openbench{} 自动选}

\begin{minted}{yaml}
project:
  grid_res: 0.5    # 目标 0.5°

reference:
  Evapotranspiration: GLEAM_v4.2a   # 自动追加 _LowRes
\end{minted}

解析规则：

\begin{itemize}
  \item \yamlkey{project.grid_res} 在 [0.5, 1.0] 区间 → 优先 \texttt{\_LowRes}
  \item \yamlkey{project.grid_res} = 0.25 → 优先 \texttt{\_MidRes}
  \item 都不存在则尝试无后缀名
\end{itemize}

\subsection{方式 2：写完整名严格指定}

\begin{minted}{yaml}
reference:
  Evapotranspiration: GLEAM_v4.2a_MidRes   # 强制 0.25°，即使 grid_res=0.5
\end{minted}

完整名永远精确匹配，不做后缀转换。

\begin{tipBox}
\textbf{什么时候该写完整名}：要求精确分辨率，或者多个 simulation 各用不同分辨率的同一数据源时。
\end{tipBox}

\section{命令：浏览 registry}

\subsection{\cli{openbench data list}}

列出所有 reference 数据集：

\begin{minted}{bash}
openbench data list
\end{minted}

\begin{exampleBox}[输出片段]
\begin{minted}{text}
Name                           Category   Type   Res      Years          Variables
─────────────────────────────────────────────────────────────────────────────────
AH4GUC_LowRes                  Urban      grid   0.5°     ?              1
BRSC_Monthly                   Water      stn    stn      ?              1
Beck_2017_Daily                Water      stn    stn      ?              1
CARE_LowRes                    Energy     grid   0.5°     ?              4
CARE_MidRes                    Energy     grid   0.25°    ?              4
CERES_EBAF_Ed4.2_LowRes        Energy     grid   0.5°     ?              7
...

Total: 98 datasets
\end{minted}
\end{exampleBox}

按变量过滤：

\begin{minted}{bash}
openbench data list --variable Evapotranspiration
\end{minted}

只列出含 \var{Evapotranspiration} 的数据集（GLEAM、FLUXCOM、ET\_Xu\_etal\_2025、FLUXNET 等）。

\subsection{\cli{openbench data path}}

查询某数据集的本地路径：

\begin{minted}{bash}
openbench data path GLEAM_v4.2a_LowRes
\end{minted}

输出：

\begin{minted}{text}
/data/openbench/refs/Grid/LowRes/Water/Evapotranspiration/GLEAM_v4.2a
\end{minted}

支持 base 名查询所有分辨率变体：

\begin{minted}{bash}
openbench data path GLEAM_v4.2a
\end{minted}

输出：

\begin{minted}{text}
GLEAM_v4.2a_LowRes: /data/openbench/refs/Grid/LowRes/...
GLEAM_v4.2a_MidRes: /data/openbench/refs/Grid/MidRes/...
\end{minted}

\section{命令：注册新数据集}

\subsection{\cli{openbench data register}}

把仓库里没有的本地数据集纳入 registry。完整签名：

\begin{minted}{bash}
openbench data register NAME --root-dir /path/to/data \
  --data-type {grid|stn} --tim-res {Month|Day|Hour|Year} \
  --grid-res 0.5 --category Water \
  --years 2000 2020 \
  -v "VarName:nc_var_name:unit"  # 可重复
\end{minted}

\begin{exampleBox}[注册一个 grid 数据集]
\begin{minted}{bash}
openbench data register MyData \
  --root-dir /data/myref \
  --data-type grid --grid-res 0.5 --tim-res Month \
  --years 2000 2020 --category Water \
  -v "Evapotranspiration:ET:mm day-1"
\end{minted}
\end{exampleBox}

\subsection{站点数据的特殊参数}

\begin{exampleBox}[注册站点数据]
\begin{minted}{bash}
openbench data register PLUMBER2 \
  --root-dir /data/PLUMBER2/dataset \
  --data-type stn --tim-res Day \
  --fulllist /data/PLUMBER2/list/PLUMBER2.csv \
  -v "Latent_Heat:Qle_cor:W/m2" \
  -v "Sensible_Heat:Qh_cor:W/m2"
\end{minted}
\end{exampleBox}

\subsection{Fallback 与单位换算}

某些数据集在不同年份/版本里同一物理量用不同变量名。用 \texttt{-f} 指定回退：

\begin{exampleBox}[fallback + 单位换算]
\begin{minted}{bash}
openbench data register ERA5 \
  --root-dir /data/era5 \
  -v "Latent_Heat:slhf:W m-2" \
  -f "Latent_Heat:surface_latent_heat_flux:J m-2:value / 3600"
\end{minted}
\end{exampleBox}

主变量是 \texttt{slhf}；如果 NC 文件里没有 \texttt{slhf} 但有 \texttt{surface\_latent\_heat\_flux}，就用回退，并应用 \texttt{value / 3600} 把焦耳/秒累积量换算为瓦特。

\subsection{交互式注册}

不带 \texttt{-v} 时进入交互模式：

\begin{minted}{bash}
openbench data register MyData --root-dir /data/myref --data-type grid
# 然后跟着提示输入 Standard variable name / NC name / unit / prefix / suffix
\end{minted}

\subsection{自定义 filter（高级）}

如果你的数据集需要 \emph{compute 表达式 / prefix-suffix} 表达不了的预处理逻辑——典型如\emph{动态生成站点列表}（按多分辨率 lookup、按 CaMA-Flood 路由约定查文件、跨变量决策选站）——可以写一个 Python filter 脚本。

\openbench{} 在加载数据时按以下顺序找 filter：

\begin{enumerate}
  \item 用户目录：\file{\$OPENBENCH_CUSTOM_DIR/<name>\_filter.py} 或 \file{\$HOME/.openbench/custom/<name>\_filter.py}
  \item 包内置：\file{<package>/data/custom/<name>\_filter.py}
\end{enumerate}

\texttt{<name>} 必须与你在 \yamlkey{reference} 里引用的数据集名一致（或去掉尾部数字后的 base 名，如 \texttt{CoLM2024} 自动回退到 \texttt{CoLM}）。模块需暴露 \texttt{filter\_<name>(info, ds=None)} 函数，分两种模式工作：

\begin{itemize}
  \item \textbf{初始化模式}（\texttt{ds=None}）：生成站点列表 CSV，写到 \texttt{info.ref\_fulllist}
  \item \textbf{运行时变量提取模式}（\texttt{ds=Dataset}）：返回 \texttt{(info, xr.DataArray)}
\end{itemize}

最简起步：

\begin{minted}{bash}
# 1. 复制内置 filter 作模板
mkdir -p ~/.openbench/custom
cp <openbench>/src/openbench/data/custom/GEBA_filter.py \
   ~/.openbench/custom/MyData_filter.py

# 2. 改函数名 filter_GEBA → filter_MyData
# 3. 改数据集路径、变量名、站点维度等
# 4. 在 reference YAML 引用 MyData
#    reference:
#      <Variable>: MyData
\end{minted}

完整契约（info 字段、CSV 列名、错误处理约定）+ 完整模板见开发卷第 4 章 §"custom: 用户可扩展 filter"。

\warninline{filter 不要调 \texttt{sys.exit} ——会杀整个 evaluator。配置错误用 \texttt{raise ValueError(...)}，runner 会捕获并把那个 (sim, ref) 标记为失败而不影响其他任务。}

\subsection{\cli{openbench data scan}：批量自动注册}

如果你已经按 \openbench{} 约定目录布局组织了一批数据集，可以一键扫描注册：

\begin{minted}{bash}
openbench data scan /Volumes/work/Reference            # 交互式
openbench data scan /Volumes/work/Reference --auto     # 全部自动注册
openbench data scan /Volumes/work/Reference --dry-run  # 预览不写
\end{minted}

\paragraph{期望的目录布局（硬编码）}

\openbench{} 严格按以下布局识别数据集：

\begin{minted}{text}
<ref_root>/
├── Grid/
│   ├── LowRes/<category>/<variable>/<dataset>/*.nc
│   ├── MidRes/<category>/<variable>/<dataset>/*.nc
│   └── HigRes/<category>/<variable>/<dataset>/*.nc
└── Station/<category>/<variable>/<dataset>/*.nc
\end{minted}

\texttt{<category>} 用 \texttt{Water} / \texttt{Heat} / \texttt{Bio} / \texttt{Meteo} / \texttt{Anth} / \texttt{Composite} 等约定名（见 \modname{registry.scanner.CATEGORY_MAP}）；\texttt{<variable>} 用标准变量名（如 \var{Evapotranspiration}）；\texttt{<dataset>} 是数据集目录名（成为 registry 名的一部分）。布局不符的目录无法自动扫描——只能用 \cli{data register} 手工注册。

\paragraph{支持的扫描深度（3 级）}

NC 文件可以位于以下任一层级：

\begin{itemize}
  \item \texttt{<dataset>/*.nc}（最常见）
  \item \texttt{<dataset>/<single\_subdir>/*.nc}（如 \texttt{0p25deg-daily/}）
  \item \texttt{<dataset>/<single>/<single>/*.nc}（如 \texttt{0p25deg/daily/}）
\end{itemize}

更深的层级会触发警告并跳过（\texttt{"beyond supported 3-level depth"}）。

\paragraph{多分支模糊跳过}

如果 \texttt{<dataset>/} 下\emph{多个}子目录都含 NC 文件（典型：composite 数据集如 \texttt{Crop/GDHY2019ver/\{maize,soybean,wheat\}/}），扫描\emph{不会}自动合并——会打 \texttt{"multiple NC-bearing subdirectories"} 警告并跳过，提示用户写 \texttt{reference\_profiles.yaml} profile 或手工注册。

\paragraph{自动从 NC 提取的字段}

\begin{itemize}
  \item \yamlkey{varname} / \yamlkey{varunit}（从 NC 第一个非坐标变量；多变量时交互式选）
  \item \yamlkey{tim_res}（从 \texttt{time} 维步长，覆盖 Hour/3Hour/6Hour/Day/8Day/Month/Year）
  \item \yamlkey{grid_res}（从 \texttt{lat} 维步长）
  \item \yamlkey{data_type}（lat 或 lon 任一 >1 → grid，否则 stn）
  \item \yamlkey{years}（从所有 NC 文件名年份正则：4 位连续数字，且\emph{不被}字母前缀如 \texttt{v2010} 污染）
  \item \yamlkey{prefix} / \yamlkey{suffix}（按文件名年份位置自动切分）
  \item \yamlkey{data_groupby}（按文件名日期粒度：Single / Day / Month / Year）
\end{itemize}

\paragraph{重扫安全性：保留的手编字段}

每次 \cli{data scan} / \cli{data register} 重扫已有数据集时，下列\emph{用户字段}从已有 catalog 中保留，不会被自动检测覆盖：

\begin{itemize}
  \item \yamlkey{description}（自由文本，扫描永远填不出有意义的）
  \item \yamlkey{category}（用户可能重新分类，如 Water → Water-Energy）
  \item \yamlkey{years}（用户对年份的修正比文件名提取更可靠）
  \item \yamlkey{timezone}（默认 0 即 UTC，但用户可改为 -8 等本地时区）
  \item \yamlkey{fulllist}（站点列表 CSV 路径——如果指向真实存在的文件就保留，不重生成）
  \item 每个变量的 \yamlkey{varname} / \yamlkey{varunit} / \yamlkey{prefix} / \yamlkey{suffix}
\end{itemize}

而 \yamlkey{tim_res} / \yamlkey{grid_res} / \yamlkey{data_type} 等"技术字段"会按新扫描结果覆盖（NC 内容才是事实源）。

\paragraph{安全保证：写入备份与并发锁}

\begin{itemize}
  \item 写入 catalog 前自动备份到 \file{<catalog>.yaml.bak}（单槽，每次写覆盖）
  \item Catalog 文件 YAML 损坏时硬错（\texttt{RuntimeError}）而非静默清空
  \item 多个 \cli{data scan} 进程并发时通过 \file{<catalog>.lock} 文件 \texttt{flock} 序列化（HPC 共享 registry 场景）
\end{itemize}

\section{命令：数据下载与状态（部分未实现）}

\subsection{\cli{openbench data download}}

\warninline{当前版本 \cli{openbench data download} \emph{尚未实现}（v3.0 alpha）。运行会得到 "Dataset download not yet implemented (requires hosted data repository)."。在线托管的数据仓库还在筹建。}

\textbf{替代方案}：reference 数据集需要自行从原始机构下载并放入约定路径。每个数据集 catalog 中的 \yamlkey{root_dir} 指明默认期望路径；用 \cli{openbench data path NAME} 查询。

\subsection{\cli{openbench data status}}

\begin{minted}{bash}
openbench data status
\end{minted}

\warninline{当前版本只报告 "Registry: N datasets available; Local cache: not yet implemented"。本地缓存检测功能在路上。}

实际可用的 fallback：

\begin{minted}{bash}
# 直接看磁盘上 reference 数据根目录
ls -la $OPENBENCH_REF_ROOT/Grid/LowRes/

# 或对单个数据集
openbench data path GLEAM_v4.2a_LowRes
ls -la "$(openbench data path GLEAM_v4.2a_LowRes)"
\end{minted}

\section{变量到数据集速查（节选）}

下面列出最常用 10 个变量的可选 reference。完整 56 变量的全部映射见附录 C：

\begin{longtable}{l p{6cm} p{4cm}}
  \toprule
  变量 & Grid 数据集 & Station 数据集 \\
  \midrule
  \endhead
  \var{Evapotranspiration} & GLEAM\_v4.2a, FLUXCOM-X-BASE, ET\_Xu\_etal\_2025（Low/Mid Res） & FLUXNET\_PLUMBER2, Xu\_et\_al2025\_stn \\
  \var{Latent\_Heat} & ERA5LAND, FLUXCOM\_LowRes & FLUXNET\_PLUMBER2, FLUX\_PLUMBER2, GEBA \\
  \var{Sensible\_Heat} & ERA5LAND, GLEAM\_v4.2a & FLUXNET\_PLUMBER2, GEBA \\
  \var{Net\_Radiation} & CLARA\_3, CERES\_EBAF, ERA5LAND, Rn-Xu2022\_MidRes & FLUXNET\_PLUMBER2, GEBA \\
  \var{Total\_Runoff} & GRUN\_ENSEMBLE, GRADES, GRFR, ERA5LAND & --- \\
  \var{Streamflow} & --- & GRDC\_Daily/Monthly, Caravan\_Daily, BRSC\_Monthly, GSIM, GSHA, ... (15+ 站点源) \\
  \var{Gross\_Primary\_Productivity} & FLUXCOM-X-BASE, GloFlux & FLUXNET\_PLUMBER2 \\
  \var{Surface\_Soil\_Moisture} & GLEAM\_v4.2a, ERA5LAND, SMAPL3 & --- \\
  \var{Precipitation} & GPCC, GPCPv3.2, IMERG\_v7, ERA5LAND, MSWEP, MSWX, CN05.1 & --- \\
  \var{Surface\_Air\_Temperature} & ERA5LAND, MSWX, CRU\_TS\_4.08, CN05.1 & --- \\
  \bottomrule
\end{longtable}

\begin{tipBox}
\var{Streamflow} 几乎全部是 station 类（grid 数据无法采集真实河流流量）。如果你的模型输出 grid 流量，\openbench{} 会按站点位置在 grid 中采样，再与站点观测对比。
\end{tipBox}

\section{选数据集的实用建议}

\subsection{优先内置数据集}

内置 registry 的数据集已经做过元信息标准化（坐标系、时间约定、单位）。新加自己的数据集（\cli{data register}）需要你保证 NetCDF 兼容；遇到坐标轴反向、时间偏移、单位不一致都会出问题。新手优先用内置。

\subsection{LowRes vs MidRes 选择}

\begin{itemize}
  \item 全球评估 + 月尺度：\texttt{\_LowRes}（0.5°）足够
  \item 区域评估 + 高分辨率比较：\texttt{\_MidRes}（0.25°）
  \item 极少需要原始（更高）分辨率，因为目标分辨率由 \yamlkey{project.grid_res} 决定，再细的源数据也会被重采样
\end{itemize}

\section{每变量多 reference：多源对比}

\openbench{} v3.0 schema：\yamlkey{reference.sources} 类型是 \texttt{dict[str, str | list[str]]}，每个变量可单 ref（字符串）或多 ref（列表）。

\subsection{为什么多源}

科学层面：单一观测产品都有偏差。Urban 与陆面温度类变量惯常双源（\texttt{["MODIST", "TRIMS\_LST"]}）；蒸散发对比常 3--4 源（\texttt{GLEAM} + \texttt{FLUXCOM-X-BASE} + \texttt{ET\_Xu\_etal\_2025}）；径流可同时对 \texttt{GRUN\_ENSEMBLE} 与 \texttt{Beck\_2017}。多 ref 评估能：

\begin{itemize}
  \item 让你看出 sim 表现是否对 ref 选择敏感（鲁棒性）
  \item 暴露 ref 之间的不一致（不同观测产品在同一区域可能差 2$\times$）
  \item 论文里同时给"vs 多个观测产品"的 metric 表，避免审稿人挑刺
\end{itemize}

\subsection{语法：三种等价形式}

\begin{minted}{yaml}
reference:
  Land_Surface_Temperature:           # list（推荐，最清晰）
    - MODIST
    - TRIMS_LST

  Evapotranspiration: "GLEAM_v4.2a_LowRes, FLUXCOM-X-BASE_LowRes"
                                       # 逗号分隔字符串（loader 自动拆）

  Total_Runoff: GRUN_ENSEMBLE_LowRes   # 单 ref 仍写字符串
\end{minted}

\subsection{(var × sim × ref) 笛卡尔积}

\openbench{} 对每个变量遍历 \emph{全部} (sim, ref) 组合做独立评估。例：

\begin{itemize}
  \item 3 个变量
  \item 4 个 simulation
  \item 变量 1 配 1 ref，变量 2 配 3 ref，变量 3 配 2 ref
  \item → 任务总数 $= 4 \times (1 + 3 + 2) = 24$ 个独立评估
\end{itemize}

每个评估：重网格 + metric/score 计算 + 出图独立产出。

\subsection{运行时间影响}

近似线性放大：m 个 sim × n 个 ref = m × n 个评估单元。但有两点优化让代价低于纯线性：

\begin{itemize}
  \item Reference 预处理（重网格、时间裁剪）\textbf{每个 ref 只做一次}，跨 sim 共享
  \item Sim 预处理\textbf{每个 sim 只做一次}，跨 ref 共享（按 v3.0 \modname{runner.local} 的 dedupe key 控制）
  \item 真正翻倍的是 evaluation 循环本身（metric 计算 + 出图）
\end{itemize}

经验：3 sim × 3 ref vs 3 sim × 1 ref 的 wallclock 比值约 2.0--2.5×，不是 3.0×。

\subsection{输出与缓存：每对独立}

输出文件命名包含 ref\_source：

\begin{itemize}
  \item Grid: \file{<var>\_ref\_<R>\_sim\_<S>\_<metric|score>.nc}
  \item Stn: \file{<var>\_stn\_<R>\_<S>\_evaluations.csv}
\end{itemize}

不同 ref 的输出绝不混淆。Cache 哈希也按 \texttt{(var, sim, ref)} 三元组建，意味着：

\begin{itemize}
  \item 改动一个 ref（换数据集、版本升级）只重算用了它的 (sim, ref) 对，不影响别的 ref
  \item 一个 (sim, ref) 失败不会阻塞其他 (sim, ref) 继续跑
  \item Comparison/统计阶段可单独看每个 ref 视角，也可跨 ref 平均（取决于你的 \yamlkey{comparison.items} 配置）
\end{itemize}

\noteinline{多 reference 是从 v2.x 沿用的能力 —— 例如 Urban 类变量惯常用 \texttt{["MODIST", "TRIMS\_LST"]} 双源。v3.0a1 早期曾因 schema 简化丢失，已修复（commit \texttt{c9fcbdc}）。}

\subsection{Provenance 警告}

\cli{openbench check} 输出中会标注 reference 元数据的可信度：

\begin{itemize}
  \item \texttt{✓}：高置信，元数据来自数据集自身或 model profile
  \item \texttt{⚠ default}：低置信，使用了默认值（可能不准）
  \item \texttt{\textasciitilde inferred}：中置信，从目录结构推断
\end{itemize}

\yamlkey{project.strict_reference: true} 把低置信度警告升级为错误，强制你显式指定。

\section{下一步}

下一章详解 \yamlkey{simulation} 块：

\begin{itemize}
  \item 21 个内置 model profile 速览
  \item 用 profile vs 自定义模型
  \item 多 simulation 对比配置
  \item time alignment 三种策略实战
  \item 单位换算与变量映射机制
\end{itemize}
